\scanpage{10}
\handwritten{24.4,67}\\

Dear Murre,\\

I am sending you enc\unsure{\deleted{o}l|\deleted{l}o}osed the sketch which I promised on ramification
% The final n of ramification is partly cut by the right paper edge; its left stem and the beginning of its upper shoulder remain visible.
% In the fourth cell of the opening word, an ascender with serifs is superposed on a rounded lower impression. The o/l correction order is not established.
data, as well as your manuscript together with my annotations. As the paragraphs 1 to 5 will have to be changed entirely, there is maybe not much point for you to look up to carefully my annotations on these. \deleted{As} I would only suggest that you put more for\unsure{mm|m}ulas, arrows etc on \deleted{\unsure{separate}} separate lines, for easier reading; that you normalize the sta\unsure{df|nf|nd}ard references to SGA, EGA (for instance, no need to indicate SGA when reference is being made to the seminar 1960-61 in which your talk is going to be included itself, reference as (XI 6.1) is then sufficient); that you be cautious when using a locution like ``if primitive $n$.th roots exist over $S$'', which does not mean much by itself, unless you make clear what you really mean. I would propose the following tentative plan for your talk:
% The right side of the o in the manuscript line ending going to is clipped, but most of its curved outline is visible.
% The crowded impression run in formulas/formmulas and the two letters after sta do not establish a recoverable correction order.

\begin{enumerate}[label={}]
\item[1.] Kummer coverings with respect to a given set of sections.
\item[2.] Ramification data.
\reversemarginpar
\margin{\figure{\begin{tikzpicture}[x=.48pt,y=.48pt,line cap=round]
\draw[gray!55,line width=.25pt] (28,64) .. controls (30,50) and (34,33) .. (30,23) .. controls (28,11) and (16,3) .. (0,0);
\end{tikzpicture}}}
\normalmarginpar
% Faint curved trace in the left margin, alongside the gap between items 2 and 3. Whether it is an intentional graphite mark or a surface trace is uncertain; no linguistic meaning is assigned.
\item[3.] Properties of the catégory of coverings of type $\underline{R}$. The group ${}_{1}^{\underline{R}}(S,\ )$
% The underlined R is above 1, with no visible base glyph. The second argument is blank. The closing parenthesis approaches the paper edge; no following full stop is visible.
\item[4.] Functorial behaviours.
\item[5.] Galois coverings of type $\underline{R}$.
\item[6.] Abhyankars theorem (for divisors with normal crossings).
\end{enumerate}
etc

The sections 6 to 10 essentially the same as in your notes, just taking into account my observations, and changing inner references according to the changes made in 1. to 5. A paragraph 11 could be an appendix on the tame variant of Lefschetz's theorem. - Par.1 is essentially your previous par.1 ; you should check exactly what is needed in it for \unsure{\deleted{f}w|\deleted{w}f}hat
% The final word contains both an f-like upper hook and w-like lower strokes; their correction order is not established.
% archive sheet number: 9
